% !TeX root = ../thuthesis-example.tex

% 中英文摘要和关键字
\begin{abstract}
    企业级大语言模型（LLM）应用在工业落地中普遍面临检索语义稀释、跨文档拓扑断裂、生成回答缺乏反思自愈能力、长事务工作流阻塞计算资源以及动态工具调用安全风险等关键挑战。针对上述技术瓶颈，本文设计并实现了一个基于 RAG（检索增强生成）与多智能体协同的企业级知识服务平台（Knowledge Hub）。

    平台创新性地提出了"控制面＋认知面"物理双核解耦架构。控制面基于 Spring Boot 3 构建，负责多租户工作区、RBAC 细粒度权限管控、知识资产元数据管理与 DAG 工作流编排；认知面以自研 Hermes 认知内核为枢纽，基于 HTTP/2 协议栈与 gRPC 双向背压流式管道实现跨进程高吞吐通信，并构建了客户端取消信号向底层大模型物理调用的即时级联中断机制。在知识检索层面，提出结构感知 Markdown AST 状态机切块算法与父子分块拓扑，实现了代码围栏原子保护、多级标题面包屑注入与表格表头跨块传播；构建了涵盖向量检索、关键词检索、元数据硬过滤与知识图谱拓扑推理的四路并发召回管线，引入倒数排名融合（RRF）与动态弱路径剪枝算法，并结合无插件应用层带偏置个性化 PageRank（PPR）实现多跳子图推理。在认知推理层面，提出了带自洽性检验与 AI Judge 三维质量评估的 Hermes 反思闭环，设计了基于规范化参数哈希指纹的 ReAct 振荡熔断机制（ReActCycleGuard）与仿生分层认知记忆体系，基于动态自适应艾宾浩斯强化衰减模型实现了记忆的多目标 Pareto 排序。在工作流编排层面，基于 Kahn 拓扑排序与 BFS 分层调度实现了节点级高并发执行与动态递归剪枝，结合不可变哈希映射压缩树（HAMT, $B=32$）与定界延续（Delimited Continuation）机制实现了人机协同审批（HITL）的零物理线程阻塞与密码学快照时间旅行。

    系统全面基于 Java 21 虚拟线程特性，严格以 DeepSeek 官方 API 作为唯一生成推理底座，以阿里百炼千问作为唯一 1536 维向量模型。实验评估表明，平台在切块内聚度、四路混合检索召回率（Recall@10 达 95.8\%，MRR 达 0.892）、幻觉抑制率（降低至 11.4\%）以及工作流并发扩展比上均显著优于现有基线方案，整体架构达到了工业级电信与金融环境所要求的鲁棒性与高可用标准。

    \thusetup{
        keywords = {检索增强生成；多智能体系统；知识图谱；Hermes 认知内核；工作流编排；定界延续；HAMT 数据结构},
    }
\end{abstract}

\begin{abstract*}
    Industrial deployment of enterprise-grade Large Language Model (LLM) applications consistently encounters critical challenges, including semantic dilution in retrieval, cross-document topological fractures, lack of self-reflective healing in generation, thread starvation in long-running workflows, and zero-trust vulnerabilities in dynamic tool calling. To overcome these fundamental bottlenecks, this paper presents Knowledge Hub, an enterprise-grade AI-native platform integrating high-fidelity Retrieval-Augmented Generation (RAG) and multi-agent cognitive orchestration.

    The platform proposes an innovative dual-core decoupled architecture comprising a Control Plane and a Cognitive Plane. The Control Plane, developed on Spring Boot 3, governs multi-tenant workspaces, fine-grained RBAC permissions, knowledge asset metadata, and DAG workflow orchestration. The Cognitive Plane, centered around the proprietary Hermes cognitive engine, communicates via an HTTP/2-based gRPC bidirectional backpressure streaming pipeline, featuring an immediate cancel-propagation mechanism that halts underlying physical LLM invocations upon client disconnection. For knowledge ingestion and retrieval, we formulate a structure-aware Markdown AST chunking algorithm coupled with parent-child hierarchical topology, enforcing code fence atomic preservation, hierarchical breadcrumb header injection, and cross-chunk table header propagation. A four-way concurrent retrieval pipeline integrating dense vectors, sparse inverted keywords, metadata hard-filtering, and Neo4j topological graph traversal is synthesized via Reciprocal Rank Fusion (RRF) with adaptive weak-path pruning, complemented by an application-level biased Personalized PageRank (PPR) for multi-hop subgraph reasoning without proprietary plugin dependencies. Within cognitive orchestration, we introduce a reflective reasoning loop governed by semantic self-consistency verification and AI Judge tri-dimensional scoring, accompanied by a canonicalized JSON-SHA256 ReActCycleGuard oscillation circuit breaker and a biomimetic hierarchical memory system driven by an adaptive Dynamic Ebbinghaus Reinforcement Decay ODE model and Pareto frontier ranking. For workflow execution, a Kahn-based layered parallel topological scheduler incorporates dynamic recursive pruning and an immutable Hash Array Mapped Trie (HAMT, branching factor $B=32$) with delimited continuations, achieving zero-thread-blocking Human-in-the-Loop (HITL) suspensions and cryptographic snapshot time-travel debugging.

    Implemented on a Java 21 virtual-thread runtime, the system strictly anchors DeepSeek official APIs as the sole generation baseline and Alibaba Qwen as the sole 1536-dimensional embedding engine. Extensive empirical evaluations demonstrate that Knowledge Hub substantially outperforms competitive baselines in chunk cohesion, multi-channel retrieval effectiveness (Recall@10 of 95.8\%, MRR of 0.892), hallucination suppression (slashed to 11.4\%), and workflow concurrency speedup, fulfilling the stringent resilience and fault-tolerance requirements of production-grade telecommunication and financial infrastructures.

    \thusetup{
        keywords* = {Retrieval-Augmented Generation; Multi-Agent Systems; Knowledge Graph; Hermes Cognitive Kernel; Workflow Orchestration; Delimited Continuation; Hash Array Mapped Trie},
    }
\end{abstract*}
